\section{Summary of the Work}

This dissertation examined whether reinforcement learning algorithms can
learn profitable trading policies that utilize market and portfolio
information from the observed states when trading assets in a portfolio.

The study designed a complete experimental system for trading the SPY
exchange-traded fund. The system uses daily market data organised into
monthly episodes, includes transaction costs and action masking, and
uses a nine-feature state representation with an optional drawdown
feature. Two reward formulations were tested: the primary wealth-change
reward and an optional risk-aware reward that penalised drawdown, along
with three learning algorithms: REINFORCE, Deep Q-learning (DQN), and
Proximal Policy Optimisation (PPO).

Experimental decisions were made on the 12-month validation dataset
(2023), while keeping the 24-month test set (2024-2025) separate and
using it only for final evaluations.

The experiments followed a fixed sequence. First, the trade size was
adjusted to ensure the agents could reach a reasonable level of exposure
relative to the buy-and-hold benchmark. The behavioural experiment then
determined whether the learned policies actually respond to the
information in their states. Once these environment, calibration and
behavioural choices had been decided, the main performance of the
learned agents was evaluated and compared against the benchmark strategy
on the unseen test period.

Additional experiments then tested whether adding drawdown to the state
changed the results, whether a risk-aware reward improved loss
management in the portfolio trade, and how the different learning
methods responded to changes in transaction costs.

\section{Conclusions Against the Research Questions}

The experiments in Chapter 5 address the four questions asked in 5.1.

\subsection{Performance}

The first question considered whether the reinforcement learning agents
could outperform the buy-and-hold benchmark on unseen data.

The results do not provide any evidence that they could. No learned
agent outperformed the simple buy-and-hold strategy on the held-out test
period. Buy-and-hold produced a mean wealth change of +\$180.25 per
month, while the best learned agent, REINFORCE, achieved +\$171.51. PPO
achieved +\$127.35 and DQN +\$74.68.

The result should not be interpreted as reinforcement learning
algorithms being incapable of learning profitable trading strategies.
All agents produced positive average wealth changes. Instead, the more
specific conclusion is that none of the tested algorithms learned a
level of trading skill that could outperform the simple investment
strategy on unseen data.

Transaction costs and the trade size also cannot explain the difference,
since they were the same throughout the experiment, providing a basis
for fair comparison. The test period, however, favoured remaining
invested, with positive market performance in 17 of the 24 months. Under
these conditions, buy-and-hold was difficult to beat as it benefits from
the rising market direction.

\subsection{Behaviour}

The second question examined whether the learned policies genuinely use
the state information provided to them.

The algorithms showed a clear difference in how they used the state. All
six DQN seeds changed their actions in response to the observed state,
while none of the twelve policy-gradient seeds from REINFORCE and PPO
did so. This changes how the performance results should be interpreted.
REINFORCE came close to the benchmark in terms of return, at +\$171.51
per month, but its policies largely shifted to buying rather than using
the market information provided.

DQN however produced policies that consistently responded to changes in
the observed state but produced a lower average return than REINFORCE.
This suggests that state dependence and financial performance are not to
be treated as equal objectives in the experimental setting. An agent can
learn to use its observations without that behaviour necessarily
yielding higher returns.

This differentiation is important because an agent can perform well for
the wrong reason. During a generally rising market, a policy that
remains invested can achieve a strong return without needing to learn
when market conditions require entering, reducing, or exiting a
position.

Looking only at return would have completely hidden the distinction
between these behaviours. The behavioural analysis is therefore an
important part of evaluating learned trading policies, as it revealed a
clear difference between the value-based and policy-gradient methods.

\subsection{State Representation}

The third question examined whether adding a drawdown feature to the
state representation improved agent performance or behaviour.

Adding drawdown to the nine-feature state produced only limited changes
using the wealth-change reward. REINFORCE\textquotesingle s mean monthly
wealth change decreased by \$4.12, while DQN\textquotesingle s increased
by \$3.13. Both changes were small compared with the variation between
their random seeds and cannot be reasonably treated as reliable
improvements.

PPO\textquotesingle s mean return increased by \$32.62, but its seed
spread also changed heavily, from \$179.80 to \$94.00. This improvement
cannot confidently be attributed to the drawdown feature itself.

Overall, the results suggest that providing additional information does
not automatically produce a better policy or improve learning unless the
objective specifies it. This is consistent with an important principle
noticed from the experiments: \textbf{information in the state is only
useful when the learning objective and optimisation process give the
agent a reason to use it}. Hence, the agent must be able to learn a
relationship between the additional information and actions that improve
the objective being optimised.

\subsection{Reward Formulation}

The fourth question considered whether adding drawdown into the reward
could encourage the agents to reduce losses and generate a more
risk-aware behaviour.

The risk-aware reward reduced drawdown mainly by reducing market
participation, not by improving understanding of market conditions. With
the selected penalty of \(\symbf{\lambda = 0.10}\),
DQN\textquotesingle s mean drawdown fell by 41\%, from 0.0140 to 0.0082.
However, its mean return also fell by 32\%, from +\$77.81 to +\$52.90,
while its exposure was approximately halved. Its Sharpe ratio fell from
0.505 to 0.416, so the reduced drawdown did not result in better
risk-adjusted performance. The reward therefore encouraged the agent to
hold less of the asset rather than to identify and avoid declining
markets. This is also reflected in March 2025, when reducing exposure
would have been particularly useful: DQN\textquotesingle s exposure
during the month was higher than its overall average, and it still
experienced 46\% of the benchmark\textquotesingle s loss.

The result illustrates an important limitation of using drawdown as a
standalone penalty. The easiest way for an agent to reduce exposure to
drawdown is to reduce its participation in the market. This behaviour
can theoretically satisfy the risk conditions of the reward without
necessarily producing a better ability to identify when the market is
likely to decline.

The penalty-weight analysis proves this interpretation. A relatively
small increase in the penalty produced a major behavioural change, with
exposure falling from 0.344 at λ = 0.10 to 0.074 at λ = 0.25. This
indicates that the reward formulation created a strong correlation
between market participation and risk avoidance.

\section{Critical Appraisal of the Project}

The most important lessons from this project were from the experimental
methodology rather than the choice of reinforcement learning algorithm
alone. Several of the most useful findings came from implementation
problems identified during development and testing.

One of the strongest aspects of the project was the verification tests.
The automated tests identified several implementation problems during
development, \textbf{three of which could have affected the
experiments' validity}. These
included issues with rolling feature calculation at the end of monthly
episodes and an implementation error in the buy-and-hold benchmark that
reduced its performance by \$38.95 per month, about 21.6\%. If this
error had remained, the dissertation\textquotesingle s main conclusion
could have been misleading, as the learned agents would have appeared to
outperform the benchmark.

The verification tests also noticed the experimental design problem in
the trade size. Using the original \(0.10C_{0}\) trade size in Chapter 3
would make the maximum attainable exposure only 0.690, compared with
0.910 for buy-and-hold. This meant that the learning agents were at a
disadvantage before training even began. Selecting \(0.25C_{0}\) allowed
a fairer comparison while still letting the agents control their
positions in smaller steps than the\(0.50C_{0}\) alternative.

The complete set of 21 verification tests passed before the final
experiments were accepted, giving a higher confidence that the reported
results reflect an appropriate experimental design rather than
implementation errors.

The behavioural analysis was another important part of the study.
Without it, the REINFORCE result of +\$171.51 could have been
interpreted as evidence that the agent had learned an intelligent,
effective trading strategy. The analysis showed this interpretation was
not justified, as the policies acted like a consistent buyer because the
market was consistently rising, completely changing the interpretation
of the result.

The project also highlighted the importance of random seed variation.
The results, particularly for PPO, varied considerably between
independent training runs. PPO had a seed spread of \$179.80 in one
configuration, which was greater than its mean return. This made it
difficult to draw conclusions based on individual training runs, as that
would have been unreliable. Reporting multiple seeds widened the
evaluation, providing a more consistent view of how each algorithm
behaved, although it increased the computational and experimental
workload; however, six seeds are still lacking for more definite
statistical inference.

Overall, the project demonstrated that careful experimental design is as
important as the choice of reinforcement learning algorithm. The main
lessons came from testing the assumptions behind the experiments, rather
than from simply comparing model returns. State dependence, action
constraints, reward design, and seed variation all affect how the final
results could be interpreted. This made the experimental process more
demanding but also produced a more reliable basis for evaluating the
learned trading policies.

\section{Limitations}

The findings of this study should be considered within the scope of the
experimental design, but there are several limitations.

\subsection{One asset and a limited market period}

The experiments use SPY as the only asset, and the held-out period from
2024 to 2025 was generally favourable to holding the asset due to rising
markets. A market with a longer or more severe decline could produce
different results for the learned agents and the buy-and-hold benchmark.
The findings cannot therefore be automatically generalised to other
assets or market conditions without further testing.

\subsection{Monthly Episode Structure}

The trading environment treated each calendar month as a separate
episode. Every episode started with the same initial capital of
\$10,000, and any remaining position was closed at the end of the month.
As a result, the results do not represent a single portfolio that
continuously compounds over the full evaluation period as a real
portfolio would.

This design made the experiments easier to control by keeping the
episode length and starting conditions consistent. However, it also
means that gains or losses from one month did not affect the starting
capital or portfolio position in the following month, which does not
represent a real-world portfolio. A longer continuous portfolio
environment could therefore lead to a different behaviour, especially
for policies that require accumulated wealth or changes in portfolio
value over time.

\subsection{Fixed Trade Size}

The agents trade fixed fractions of the initial capital rather than
choosing their position size continuously. This was calibrated to create
a controlled, reasonable environment for comparison with the benchmark,
but it restricts how much the agents can express their decisions. A
continuous position-sizing model could potentially respond differently
based on its exposure in a particular market condition. Rather than
repeatedly applying fixed buy or sell actions, such an agent could
select its desired portfolio exposure, producing different results in
returns and trading behaviour.

\subsection{Limited Number of Seeds and Lack of Significance Testing}

Six seeds were used for each configuration, and the seed-level results
and variation between them are reported throughout the results. This is
stronger than relying on a single training run, but it does not provide
sufficient statistics for formal claims about small differences. For
example, the \$3.13 change observed for DQN and \textbf{−\$4.12} for
REINFORCE when the state representation was changed is not large enough
to be treated as a meaningful improvement.

\subsection{Risk Penalty Selection was based on DQN}

The reward\textquotesingle s penalty weight λ was selected using DQN
because it was the algorithm that consistently showed state-dependent
behaviour. This does not guarantee that the same penalty value is
appropriate for REINFORCE or PPO, as it was not independently optimised
for every algorithm. A separate calibration for each algorithm could
produce different results. This limitation is particularly relevant
because the algorithms responded differently to the state and reward
formulation.

\subsection{The risk penalty only targets drawdown}

The risk-aware reward penalises drawdown but does not directly measure
whether the return achieved is sufficient for the level of risk taken.
While drawdown is directly relevant to losses, it does not fully
represent the portfolio risk. A reward that relates return to risk,
rather than penalizing risk alone, was not tested.

Additionally, as shown in Section 5.7, the agents could reduce the
penalty simply by reducing their exposure to the asset or not trading at
all. The experiments showed that a lower drawdown value should not
automatically be associated with evidence of better risk management.

\section{Future Work}

The limitations above suggest several directions for furthering the
research. Each stems directly from a limitation or an unresolved result
in the experiments. The behavioural analysis should also be applied to
these future works to determine the agent's behaviour to changes
alongside performance improvements.

\subsection{Longer and More Challenging Market Periods}

The experiments could be repeated over longer market periods containing
more market declines. A useful change would be to use a walk-forward
training approach, allowing the agents to be retrained as market
conditions change. The current study provides useful evidence for a
generally rising test period, but testing market declines would provide
a balanced evaluation of the portfolio.

This would better reflect the changing nature of financial markets,
allowing the system to encounter different combinations of rising,
falling, and volatile conditions. Such an evaluation would also better
test the agents' risk-aware behaviour during downtrends.

\subsection{Predictive and Uncertainty-Aware State Information}

The drawdown feature describes the portfolio\textquotesingle s previous
losses but does not provide information about future market movements.
Future work could therefore include predictive information, such as
probabilistic forecasts containing estimates of forecast uncertainty.

This would give a state-dependent agent more useful information for
making decisions about future market conditions rather than relying only
on historical data. It would also directly test whether
DQN\textquotesingle s ability to use the state leads to more useful
market understanding and trading decisions.

\subsection{Risk-Efficiency Objectives}

The most direct correction to the drawdown penalty could be to replace
it with a reward that considers return relative to risk, such as a
Sharpe-based or drawdown-relative objective~\cite{moody2001}. This would make simply
reducing market exposure a less effective way of improving the reward,
as seen in the current experiments.

The results in Section 5.7 make this future rework highly relevant
because the tested penalty mainly caused agents to reduce exposure
rather than learn when downside risk is increasing.

\subsection{More Seeds for Better Statistical Analysis}

More random training seeds would improve the reliability of comparisons~\cite{henderson2018}
between algorithms and configurations\textbf{.} Future work should
combine this with appropriate statistical analysis to determine whether
observed differences between configurations are likely to represent
genuine results rather than random noise.

\subsection{More flexible trade sizes}

Future experiments could allow the agent to choose its position size
continuously rather than using fixed trade amounts. This would remove
the exposure limit created by the current environment and allow the
agent to express different levels of confidence by adjusting its
position size, potentially enabling smoother portfolio adjustments.

\section{Closing Remarks}

This dissertation examined whether standard reinforcement learning could
beat the simple buy-and-hold strategy on unseen data. The experiments
found no evidence that it could. The best-performing learned agent
finished \$8.74 per month below buy-and-hold, and it did so by ignoring its market
observations to follow a simple trading behaviour.

The main contribution of the study is therefore not a successful trading
strategy, but a careful examination of why the learned agents performed
as they did. The comparisons were conducted in a controlled environment
in which trade sizing was calibrated, transaction costs were applied,
and the final results were produced in a suitable experimental
framework. The policies were then evaluated on both behavioral and
financial performance, and the main findings were tested across the
different state representations and reward functions.

The implementation, experimental framework and behavioural analysis
developed in this study provide a useful basis for future work. More
advanced models can build on this framework while being evaluated and
compared against the same fair standards.
